\section{Two modules that no star of bounded depth distinguishes}\label{section: local}

\subsection{Spectra}

Let $\M$ be a valuation whose universe is an interval $\range{u, v}$ with $u < v$.
For example, $\M$ can be a module.
We call the vertices $x$ with $u < x < v$ \emph{interior}.
For an interior $x$, let
\[
  \lambda_q(x) \defeq \tp{q}(\M\range{u, x}, u, x),
  \qquad
  \rho_q(x) \defeq \set{(F, C) \in \Rep{q}^{2} \mid (F^{*} \compo (C \inter \Dt))(x, v)}.
\]
The first component describes the prefix up to $x$, and the second describes which stars can leave $x$ towards $v$.
The \emph{spectrum} of $\M$ is $\Spec{q}(\M) \defeq \set{(\lambda_q(x), \rho_q(x)) \mid x \text{ interior}}$.

\begin{definition}\label{definition: local equivalence}
  We write $\M \eqq{q} \M'$ if $\tp{q}(\M, u, v) = \tp{q}(\M', u', v')$ and $\Spec{q}(\M) = \Spec{q}(\M')$.
  Here $u, v$ and $u', v'$ are the endpoints of $\M$ and $\M'$.
\end{definition}

\begin{lemma}[Spectrum]\label{lemma: spectrum}
  Let $\M \eqq{q} \M'$.
  Then for all $G, A, F, C \in \SF{q}$, each of the following holds in $\M$ iff it holds in $\M'$:
  \[
    G(u, v) \qquad\text{and}\qquad ((A \inter \Dt) \compo F^{*} \compo (C \inter \Dt))(u, v).
  \]
\end{lemma}
\begin{proof}
  The first claim follows from \Cref{lemma: SF types}.
  For the second, we may assume $A, F, C \in \Rep{q}$.
  The observation holds iff there is $x$ with $u < x$, $A(u, x)$ and $(F^{*} \compo (C \inter \Dt))(x, v)$.
  The last condition forces $x < v$, so $x$ is interior.
  By \Cref{lemma: locality,lemma: SF types}, $A(u, x)$ depends only on $\lambda_q(x)$.
  So the observation holds iff some $(\lambda, \rho) \in \Spec{q}(\M)$ has $\lambda \models A$ and $(F, C) \in \rho$.
\end{proof}

\subsection{Prefix words and suffix words}

Let $\M = \M(w)$ be a module.
We label each interior vertex $z$ with the letter $(f(z), s(z)) \in \Gamma \defeq \set{0,1}^{2}$.
Here $f(z) = 1$ iff $\vr(u, z)$, that is, $z$ is the first vertex of a branch.
And $s(z) = 1$ iff $\vc(u, z)$, that is, $z$ is selected.
For an interior $x$, the \emph{prefix word} $\alpha_x \in \Gamma^{+}$ is the sequence of labels of the vertices in $(u, x]$.
The \emph{suffix word} $\beta_x \in \set{0,1}^{+}$ is the sequence of $f$-labels of the vertices in $[x, v)$.
Let $\alpha(\M)$ be the label sequence of all interior vertices.

\begin{lemma}[Prefixes]\label{lemma: prefixes}
  Let $\M, \M'$ be modules.
  \begin{enumerate}
    \item If $\tau_q(\alpha(\M)) = \tau_q(\alpha(\M'))$, then $\tp{q}(\M, u, v) = \tp{q}(\M', u', v')$.
    \item If $x, x'$ are interior vertices of $\M, \M'$ with $\tau_q(\alpha_x) = \tau_q(\alpha'_{x'})$, then $\lambda_q(x) = \lambda'_q(x')$.
  \end{enumerate}
\end{lemma}
\begin{proof}
  We can rebuild $\M\range{u, x}$ from the word $\alpha_x$.
  Add a new minimum $u$ to the word.
  Let $\vr$ consist of the pairs $(u, z)$ with $f(z) = 1$ and the successor pairs $(z, z')$ with $f(z') = 0$.
  Let $\vc$ consist of the pairs $(u, z)$ with $s(z) = 1$, and let $\vb$ be empty.
  This is correct because two consecutive interior vertices lie in the same branch iff the second one is not the first vertex of a branch.
  For the whole module we also add a new maximum $v$ and set $\vb \defeq \set{(u, v)}$.

  Now take a winning strategy of Duplicator in the $q$-round game on $\alpha_x$ and $\alpha'_{x'}$.
  Extend it by matching $u$ with $u'$ (and $v$ with $v'$).
  The result is a winning strategy for the rebuilt structures, with $x$ and $x'$ matched as the last positions.
  Indeed, all relations of the rebuilt structures are quantifier-free definable from the order, the successor, the letters, $u$ and $v$.
\end{proof}

\begin{lemma}[Suffixes]\label{lemma: suffixes}
  There are a finite monoid $\Qm$ and a morphism $\chi \colon \set{0,1}^{*} \to \Qm$ with the following property.
  For all modules $\M, \M'$ and interior vertices $x, x'$, if $\chi(\beta_x) = \chi(\beta'_{x'})$, then $\rho_q(x) = \rho'_q(x')$.
\end{lemma}
\begin{proof}
  The suffix $\M\range{x, v}$ contains neither $u$ nor any $\vc$- or $\vb$-pair.
  It is determined by $\beta_x$.
  Its $\vr$-pairs are the successor pairs $(z, z')$ with $f(z') = 0$ and $z' \ne v$.
  Fix $(F, C) \in \Rep{q}^{2}$, and let $L_{F, C}$ be the set of $\beta \in \set{0,1}^{+}$ for which the structure determined by $\beta$ satisfies $(F^{*} \compo (C \inter \Dt))(\min, \max)$.
  Encode the structure as the word $\beta\sharp$, where $\sharp$ stands for $v$.
  All relations are first-order definable in $\beta\sharp$, and reflexive transitive closure is definable in monadic second-order logic.
  So $L_{F, C}$ is regular by the theorem of B\"uchi, Elgot and Trakhtenbrot~\cite{buchiWeakSecondorderArithmetic1960}.
  Let $\chi$ be the product of the syntactic morphisms of the finitely many languages $L_{F, C}$.
\end{proof}

\subsection{Codes of branches}

For a branch type $t$ with length $\ell$ and selected index $\sigma$, the \emph{prefix code} $\pc(t) \in \Gamma^{\ell}$ has the following letters.
The first letter is $(1, 0)$, the $\sigma$-th letter is $(0, 1)$, and all other letters are $(0, 0)$.
The \emph{suffix code} is $\scode(t) \defeq 1\,0^{\ell - 1}$.
We extend both codes to words over $\set{\bg, \bs, \bl}$ by concatenation.
So $\alpha(\M(w)) = \pc(w)$.
For an interior $x$ of $\M(w)$, $\alpha_x$ is a prefix of $\pc(w)$ and $\beta_x$ is a suffix of $\scode(w)$.
We define the morphisms
\[
  \varphi \defeq \tau_q \circ \pc \colon \set{\bg, \bs, \bl}^{*} \to \Types{q}(\Gamma),
  \qquad
  \hat\chi \defeq \chi \circ \scode \colon \set{\bg, \bs, \bl}^{*} \to \Qm.
\]
From now on, fix $q$ and let $m \defeq 2^{q+3}$.

\begin{lemma}\label{lemma: branch identities}
  \leavevmode
  \begin{enumerate}
    \item\label{item: long} $\varphi(\bg) = \varphi(\bl)$.
          Moreover, the prefix of length $m-1$ of $\pc(\bg)$ and the prefix of length $m$ of $\pc(\bl)$ have the same $\tau_q$.
    \item\label{item: shift} $\hat\chi(\bg) = \hat\chi(\bs)$.
  \end{enumerate}
\end{lemma}
\begin{proof}
  (\ref{item: long})
  We have $\pc(\bg) = (1,0)\,(0,0)^{m-2}\,(0,1)$ and $\pc(\bl) = (1,0)\,(0,0)^{m-1}\,(0,1)$.
  The two prefixes are $(1,0)\,(0,0)^{m-2}$ and $(1,0)\,(0,0)^{m-1}$.
  Since $m - 2 > 2^{q+1}$, \Cref{lemma: word types}(\ref{item: powers}) gives $\tau_q((0,0)^{m-2}) = \tau_q((0,0)^{m-1})$.
  We conclude by \Cref{lemma: word types}(\ref{item: composition}).
  (\ref{item: shift}) $\scode(\bg) = \scode(\bs) = 1\,0^{m-1}$.
\end{proof}

So the prefix monoid cannot tell $\bg$ from $\bl$, and the suffix monoid cannot tell $\bg$ from $\bs$.
Neither monoid alone separates $\bg$ from $\Bad$.
The question is whether both monoids can be fooled \emph{at the same point}.

\subsection{The local pair}

\begin{proposition}[Local pair]\label{proposition: local pair}
  For every $q$ there are modules $\M^{+}_{q}$ and $\M^{-}_{q}$ with the following properties.
  \begin{enumerate}
    \item They have the same vertices and the same $\vr$ and $\vb$.
          Their $\vc$ differ in exactly one pair: $\M^{-}_q$ has $(u, y_{m-1})$ where $\M^{+}_q$ has $(u, y_m)$, for one branch $y_1 < \dots < y_m$.
    \item $\Kt(u, v)$ holds in $\M^{+}_q$ and fails in $\M^{-}_q$.
    \item $\M^{+}_q \eqq{q} \M^{-}_q$.
  \end{enumerate}
\end{proposition}

\begin{proof}
  Let $\Pm \defeq \varphi(\Bad^{*})$ and $\eta \defeq (\varphi, \hat\chi) \colon \Bad^{*} \to \Types{q}(\Gamma) \times \Qm$, and let $\Tm \defeq \eta(\Bad^{*})$.
  The monoid $\Pm$ is aperiodic by \Cref{lemma: word types}(\ref{item: powers}), and the first projection maps $\Tm$ onto $\Pm$.
  Fix an idempotent $e = (e_P, e_Q) \in \Kid(\Tm)$.
  By \Cref{lemma: minimal ideal},
  \begin{equation}\label{equation: collapse}
    e\Tm e \text{ is a group}, \qquad e_P\, n\, e_P = e_P \quad (n \in \Pm).
  \end{equation}
  Note that $\varphi(\bg) = \varphi(\bl) \in \Pm$ by \Cref{lemma: branch identities}.

  \emph{Reset words.}
  Choose a nonempty $t \in \Bad^{*}$ with $\eta(t) = e$.
  If $e$ is not the identity, any preimage works.
  Otherwise $\Tm = \Kid(\Tm)$ is a group, and we take $t \defeq \bs^{\card{\Tm}}$.
  For $x \in \Bad$, the element $e\eta(x)e$ lies in the group $e\Tm e$.
  Choose $v_x \in \Bad^{*}$ such that $\eta(v_x)$ is its inverse in this group.
  Let
  \[
    L_x \defeq t\, v_x\, t\, x\, t \quad (x \in \Bad),
    \qquad
    L^{+} \defeq t\, v_{\bs}\, t\, \bg\, t.
  \]
  Then $\eta(L_x) = \eta(v_x) \cdot e\eta(x)e = e$.
  By \eqref{equation: collapse} and \Cref{lemma: branch identities}(\ref{item: shift}), also $\eta(L^{+}) = e$.
  Define
  \[
    W^{-} \defeq L_{\bs}\, L_{\bl}\, L_{\bs}\, t,
    \qquad
    W^{+} \defeq L^{+}\, L_{\bl}\, L_{\bs}\, t,
    \qquad
    \M^{\pm}_q \defeq \M(W^{\pm}).
  \]
  The two words differ only in the displayed letter of the first block, which is $\bs$ in $W^{-}$ and $\bg$ in $W^{+}$.
  Since $\bs$ and $\bg$ have the same length and differ only in the selected vertex, property~1 holds.
  Property~2 follows from \Cref{lemma: signal}, since $W^{-} \in \Bad^{*}$.

  \emph{Root types.}
  We have $\varphi(W^{\pm}) = e_P^{4} = e_P$.
  So the root types agree by \Cref{lemma: prefixes}.

  \emph{Critical zones.}
  Each of the three blocks has the form $t\, v\, t\, c\, t$.
  We call the branch coming from the letter $c$ the \emph{displayed} branch of the block.
  The \emph{critical zone} of the block consists of the last two vertices of the displayed branch and the vertices of the last copy of $t$ in the block.
  Let $Z, Z', Z''$ be the critical zones of the first, second and third blocks.
  These sets are the same in $\M^{+}$ and $\M^{-}$, since the two modules have the same vertices.
  Let $x$ be a vertex in the critical zone of a block whose displayed letter is $c$.
  Then its prefix word and suffix word factor as
  \[
    \alpha_x = \pi\, \zeta_x,
    \qquad
    \beta_x = \xi_x\, \kappa.
  \]
  Here $\pi$ is the prefix code of everything before the displayed branch, and $\zeta_x$ is the prefix of $\pc(c\, t)$ that ends at $x$.
  Further, $\xi_x$ is the suffix of $\scode(c\, t)$ that starts at $x$, and $\kappa$ is the suffix code of everything after the block.
  The word before the displayed branch is $t\, v\, t$, possibly preceded by earlier blocks.
  The word after the block is a nonempty product of blocks and $t$.
  So \eqref{equation: collapse} and $\eta(t) = \eta(L_x) = \eta(L^{+}) = e$ give
  \begin{equation}\label{equation: zone states}
    \tau_q(\pi) = e_P \quad\text{and}\quad \chi(\kappa) = e_Q
    \qquad\text{in all three blocks, in both modules.}
  \end{equation}
  Hence $\tau_q(\alpha_x) = e_P\, \tau_q(\zeta_x)$ and $\chi(\beta_x) = \chi(\xi_x)\, e_Q$.
  The point is that only $\zeta_x$ and $\xi_x$ depend on the block.
  We match the vertices of two critical zones position by position.
  The $j$-th last vertex of one displayed branch ($j = 1, 2$) goes to the $j$-th last vertex of the other, and the vertices of the final copies of $t$ correspond by position.
  We write $\lambda^{\pm}, \rho^{\pm}$ for the maps of $\M^{\pm}_q$.
  The matching is shown in \Cref{figure: reset}.

  \begin{figure}[t]
    \centering
    \begin{tikzpicture}[x=1cm,y=1cm,
    cellbox/.style={draw,minimum height=0.55cm,inner sep=0pt,font=\scriptsize},
    zone/.style={fill=orange!25},
    match/.style={-{Stealth[length=5pt]},thick}]
  % a block t v t c t starting at x=#1, row y=#2, displayed letter #3, name prefix #4
  \newcommand{\block}[4]{%
    \begin{scope}[shift={(#1,#2)}]
      \fill[zone] (2.35,-0.275) rectangle (3.45,0.275);
      \draw (0,-0.275) rectangle (3.45,0.275);
      \foreach \a in {0.55,1.35,1.9,2.85}{\draw[gray] (\a,-0.275) -- (\a,0.275);}
      \node[font=\scriptsize] at (0.275,0) {$t$};
      \node[font=\scriptsize] at (0.95,0) {$v_{\bs}$};
      \node[font=\scriptsize] at (1.625,0) {$t$};
      \node[font=\small] at (2.375,0) {#3};
      \node[font=\scriptsize] at (3.15,0) {$t$};
      \coordinate (#4) at (2.9,0);
    \end{scope}}
  \newcommand{\blockl}[4]{%
    \begin{scope}[shift={(#1,#2)}]
      \fill[zone] (2.35,-0.275) rectangle (3.45,0.275);
      \draw (0,-0.275) rectangle (3.45,0.275);
      \foreach \a in {0.55,1.35,1.9,2.85}{\draw[gray] (\a,-0.275) -- (\a,0.275);}
      \node[font=\scriptsize] at (0.275,0) {$t$};
      \node[font=\scriptsize] at (0.95,0) {$v_{\bl}$};
      \node[font=\scriptsize] at (1.625,0) {$t$};
      \node[font=\small] at (2.375,0) {#3};
      \node[font=\scriptsize] at (3.15,0) {$t$};
      \coordinate (#4) at (2.9,0);
    \end{scope}}
  % rows
  \node[anchor=east] at (-0.15,2) {$W^{+}$};
  \node[anchor=east] at (-0.15,0) {$W^{-}$};
  % W+ : L+ L_l L_s t
  \block{0}{2}{$\bg$}{P1}
  \blockl{3.6}{2}{$\bl$}{P2}
  \block{7.2}{2}{$\bs$}{P3}
  \draw (10.8,1.725) rectangle (11.35,2.275); \node[font=\scriptsize] at (11.075,2) {$t$};
  % W- : L_s L_l L_s t
  \block{0}{0}{$\bs$}{M1}
  \blockl{3.6}{0}{$\bl$}{M2}
  \block{7.2}{0}{$\bs$}{M3}
  \draw (10.8,-0.275) rectangle (11.35,0.275); \node[font=\scriptsize] at (11.075,0) {$t$};
  % block labels
  \foreach \x/\l in {1.725/{$L^{+}$},5.325/{$L_{\bl}$},8.925/{$L_{\bs}$}}{\node[font=\scriptsize] at (\x,2.5) {\l};}
  \foreach \x/\l in {1.725/{$L_{\bs}$},5.325/{$L_{\bl}$},8.925/{$L_{\bs}$}}{\node[font=\scriptsize] at (\x,-0.5) {\l};}
  % zone names
  \node[font=\scriptsize] at (2.9,1.45) {$Z$};
  \node[font=\scriptsize] at (6.5,1.45) {$Z'$};
  \node[font=\scriptsize] at (10.1,1.45) {$Z''$};
  \node[font=\scriptsize] at (2.9,0.55) {$Z$};
  \node[font=\scriptsize] at (6.5,0.55) {$Z'$};
  \node[font=\scriptsize] at (10.1,0.55) {$Z''$};
  % matchings
  \draw[match,blue!70!black] ($(P1)+(0.15,-0.3)$) -- node[pos=0.35,above,sloped,font=\scriptsize]{$\bg \mapsto \bl$} ($(M2)+(-0.1,0.3)$);
  \draw[match,red!70!black] ($(M1)+(0.15,0.3)$) -- node[pos=0.2,below,sloped,font=\scriptsize]{$\bs \mapsto \bs$} ($(P3)+(-0.1,-0.3)$);
\end{tikzpicture}

    \caption{The words $W^{-}$ and $W^{+}$ with the critical zones $Z, Z', Z''$ (shaded).
      Before each displayed branch the prefix state is $e_P$, and after each block the suffix state is $e_Q$.
      Outside $Z$ every vertex is matched to itself.
      The vertices of $Z$ in $\M^{+}_q$ (good branch $\bg$) are matched to $Z'$ in $\M^{-}_q$ (long branch $\bl$).
      The vertices of $Z$ in $\M^{-}_q$ are matched to $Z''$ in $\M^{+}_q$ (both shifted branches $\bs$).}
    \label{figure: reset}
  \end{figure}

  \begin{claim}\label{claim: outside}
    For every interior $x \notin Z$: $(\lambda^{+}(x), \rho^{+}(x)) = (\lambda^{-}(x), \rho^{-}(x))$.
  \end{claim}
  \begin{claimproof}
    The suffixes $\M^{\pm}_q\range{x, v}$ coincide, because they contain no $\vc$-pair.
    So $\rho^{+}(x) = \rho^{-}(x)$.
    If $x$ is before $Z$, both changed $\vc$-targets lie after $x$, so $\M^{+}_q\range{u, x} = \M^{-}_q\range{u, x}$.
    If $x$ is after $Z$, then $\alpha^{\pm}_x = \pc(L^{\pm})\, \gamma$ with $L^{-} \defeq L_{\bs}$ and a common word $\gamma$.
    Since $\varphi(L^{+}) = \varphi(L_{\bs}) = e_P$, \Cref{lemma: prefixes} gives $\lambda^{+}(x) = \lambda^{-}(x)$.
  \end{claimproof}

  \begin{claim}\label{claim: good to long}
    Let $x \in Z$ and let $x' \in Z'$ be its match.
    Then $(\lambda^{+}(x), \rho^{+}(x)) = (\lambda^{-}(x'), \rho^{-}(x'))$.
  \end{claim}
  \begin{claimproof}
    The displayed letters are $\bg$ (block one of $W^{+}$) and $\bl$ (block two of $W^{-}$).
    The words $\zeta_x$ and $\zeta_{x'}$ are one of the following pairs.
    (a) The prefixes of lengths $m-1$ of $\pc(\bg)$ and $m$ of $\pc(\bl)$.
    (b) The words $\pc(\bg)$ and $\pc(\bl)$.
    (c) The words $\pc(\bg)\,\delta$ and $\pc(\bl)\,\delta$ for a common prefix $\delta$ of $\pc(t)$.
    In each case $\tau_q(\zeta_x) = \tau_q(\zeta_{x'})$ by \Cref{lemma: branch identities}(\ref{item: long}) and \Cref{lemma: word types}(\ref{item: composition}).
    The words $\xi_x$ and $\xi_{x'}$ are equal.
    They are $00\,\scode(t)$ or $0\,\scode(t)$ in the branch, since $m \ge 3$ and the last two vertices are not first vertices of the branch.
    In the copy of $t$ they are the same suffix of $\scode(t)$.
    By \eqref{equation: zone states} and \Cref{lemma: prefixes,lemma: suffixes}, the claim follows.
  \end{claimproof}

  \begin{claim}\label{claim: shifted to shifted}
    Let $x \in Z$ and let $x'' \in Z''$ be its match.
    Then $(\lambda^{-}(x), \rho^{-}(x)) = (\lambda^{+}(x''), \rho^{+}(x''))$.
  \end{claim}
  \begin{claimproof}
    Both displayed letters are $\bs$.
    Hence $\zeta_x = \zeta_{x''}$ and $\xi_x = \xi_{x''}$, and we conclude by \eqref{equation: zone states} as before.
  \end{claimproof}

  By \Cref{claim: outside,claim: good to long}, $\Spec{q}(\M^{+}_q) \subseteq \Spec{q}(\M^{-}_q)$.
  By \Cref{claim: outside,claim: shifted to shifted}, $\Spec{q}(\M^{-}_q) \subseteq \Spec{q}(\M^{+}_q)$.
  Together with the root types, $\M^{+}_q \eqq{q} \M^{-}_q$.
\end{proof}

\begin{remark}
  The term $\Kt$ itself has star height one and distinguishes $\M^{+}_q$ from $\M^{-}_q$.
  So $\eqq{q}$ does not imply equivalence for all terms of star height one.
  What $\eqq{q}$ does guarantee is that stars of depth at most $q$ \emph{passing through} a module do not see the difference (\Cref{lemma: replacement}).
  The amplification in \Cref{section: cells} makes this enough.
  The suffix monoid $\Qm$ may contain nontrivial groups, for instance because $(\vr \compo \vr)^{*}$ counts modulo two.
  This is why we need the minimal ideal and the inverses $v_x$, and not only the aperiodicity of the prefix monoid.
\end{remark}
